% Cote 119: offprints, reports, courses, CV and letters (typed and printed part): archive cote 119, site page 228 (PDF page 229).
% typescript. Draft compared: none.
% AI-transcribed by Claude Opus 5.5 (claude-opus-5-5, adaptive thinking, effort high), Message Batches API, batch r119_full (msgbatch_01DwZNm5Y7YsKBWyAN4sFs8w), 2026-11-13; not yet checked by eye.
\scanpage{119/228}
\mspage{. 15 .}

du groupe fondamental dans les premiers cas non triviaux,
comme celui d'une courbe algébrique privée de certains
points ; on peut y adjoindre les théorèmes de génération
et de présentation finie du groupe fondamental d'une
variété algébrique sur un corps algébriquement clos.
Ce programme est accompli pour l'essentiel dans SGA 1,
en utilisant à la fois les résultats classiques sur le
corps des complexes (établis par voie transcendante)
et une panoplie d'outils faits sur mesure (théorie de
la descente, étude des morphismes étales, théorèmes
d'existence de faisceaux cohérents ...). Les autres
références contiennent des résultats plus spéciaux :
théorèmes du type de Lefschetz dans SGA 2, action des
groupes de monodromie locale sur le groupe fondamental
d'une fibre dans SGA 7, Exp. I, calculs de certains
groupes fondamentaux locaux dans [19 bis], via les groupes
fondamentaux de certains schémas formels. Tous ces
résultats ont été utilisés couramment dans de nombreux
travaux, et en ont inspiré d'autres comme la thèse de
Mme Michèle Raynaud [\unsure{$\ast$}].

f) \uline{Théorème de Lefschetz locaux et globaux} pour les
groupes de Picard, le groupe fondamental, la cohomologie
étale, la cohomologie cohérente. Il s'agit ici de la
comparaison entre les invariants (cohomologiques ou
homotopiques) d'une variété algébrique et d'une section
hyperplane. Les idées de départ sont développées dans
[16, SGA 2]. Cependant, pour des énoncés ``définitifs'',
en termes de conditions nécessaires et suffisantes, se
reporter plutôt à la thèse de Mme Michèle Raynaud [\unsure{$\ast$}]
déjà citée.

\mspage{\hand{228}}
